\par\SourceLeaf{492}%
Im bevorstehenden Sommer wird derselbe zuerst auf der Nordspitze von Jütland die astronomischen Beobachtungen anstellen, und nachher theils in Lauenburg mit neuen Instrumenten die vorjährigen Messungen wiederholen und die Azimuthalbestimmungen machen, theils bei Hamburg eine Grundlinie, einige Meilen lang, messen. Nur für einen Theil der Operationen in Lauenburg und die Basismessung, nicht aber für die Operationen auf der Nordspitze von Jütland, wird meine Gegenwart gewünscht. Das Dänische Gouvernement, indem es meine Anwesenheit bei diesen delicaten Operationen gewünscht hat, scheint mir von der Ansicht ausgegangen zu sein, dass vereinte Berathung bei diesem lediglich zum Besten der Wissenschaft unternommenen Geschäft, förderlich sein werde, was durch blossen Briefwechsel der Natur der Sache nach nur unvollkommen oder gar nicht erreicht werden könne. Zweitens, da in Zukunft die ganze Gradmessung in extenso in einem eignen Werke der gelehrten Welt bekannt gemacht werden soll, wird die Anwesenheit eines Zeugen bei einigen der wichtigsten Operationen dazu dienen, die Authenticität zu verstärken. So wurden im Jahre 1798 zu der Basismessung bei Paris alle damals mit Frankreich nicht im Krieg begriffene Staaten von Frankreich eingeladen, qualificirte Astronomen abzuschicken, was auch von Dänemark, der Schweiz, Holland, Spanien und andern Staaten geschah. Diese beiden Rücksichten beziehen sich zunächst nur auf die Dänischen Gradmessungen für sich allein genommen. In der Voraussetzung, dass früh oder spät eine weitere Fortsetzung durch das Hannöversche angeordnet werden könnte, scheint es aber noch besonders wichtig zu sein, dass durch einen diesseitigen Astronomen das Vertrauen gehörig gewürdigt werden könne, welches diejenigen Operationen besonders verdienen, an welche eine solche Fortsetzung unmittelbar sich anschliessen müsste.

Nach diesen Betrachtungen bin ich sehr gern bereit, der mir schmeichelhaften Aufforderung Folge zu leisten, und werde ich den Dänischen Astronomen ersuchen, mich von der Zeit, wo jene Operationen werden vorgenommen werden, so früh als möglich zu benachrichtigen, damit ich mich bei meinen anderweitigen Geschäften danach einrichten könne. Mit den Vorlesungen wird dies wol keine Schwierigkeiten haben, da doch jene Operationen erst im Spätsommer anfangen werden. Leid wird es mir zwar thun, aller Wahrscheinlichkeit nach die in ihrer Art einzige am 7. Sept. eintreffende grosse ringförmige Sonnenfinsterniss hier nicht mehr beobachten zu können; allein ich werde mich damit beruhigen, dass theils dieselbe hier doch durch den Prof. \textsc{Harding} wird beobachtet werden können, theils, dass dieselbe auch im Holsteinischen gleichfalls ringförmig erscheinen und es leicht thunlich sein wird, dort gute, wenn auch den hiesigen nicht ganz gleich kommende Hülfsmittel zur Beobachtung dieses merkwürdigen Phänomens herbeizuschaffen. — — —

Göttingen den 27. Februar 1820.
\begin{flushright}C. F. \textsc{Gauss}.\end{flushright}

\par\bigskip\Needspace{8\baselineskip}
\pdfbookmark[0]{P. M. betreffend die Hannoversche Gradmessung.}{pm-hannoversche-gradmessung}
\begin{center}\large P. M. BETREFFEND DIE HANNOVERSCHE GRADMESSUNG.\end{center}
\par
— — — Ich habe bereits in meinem frühern Bericht das für dieses Geschäft unentbehrliche Bedürfniss mehrerer Hauptinstrumente auseinandergesetzt. Solange es sogar noch ungewiss war, woher dieselben zu erhalten stehen würden, konnte natürlich mit den Operationen selbst noch kein Anfang gemacht werden, einige astronomische Vorarbeiten auf hiesiger Sternwarte, womit ich schon im vorigen Sommer angefangen, abgerechnet. Inzwischen habe ich jene Hauptbedürfnisse, namentlich einen grössern möglichst vollkommenen Theodolithen und ein sogenanntes Universalinstrument bei dem ersten jetzt lebenden Künstler, von \textsc{Reichenbach} in München, bestellt, welcher nicht nur diese Bestellung angenommen, sondern auch diese Instrumente resp. Anfang Mai und im Juli d. J. zu liefern versprochen hat.
